% GENERATED FILE. Edit canonical lessons, then run npm run generate:books.
% canonical-source-hash: fnv1a64:a6ed2b9e57557c06
% canonical-lessons: RU-C87-straight-on, RU-C87-nearby, RU-C87-airport, RU-C87-stop, RU-C87-square

\chapter{Straight on, Nearby, Airport, Stop, Square}
\label{ch:ru-straight-on-nearby-airport-stop-square}
\hlchaptermodality{87}

\begin{chapteropening}
\textbf{By the end of this chapter:} \emph{I can say straight on, nearby, an airport, a stop and a square.}

Complete the last lesson of chapter 87: I can say straight on, nearby, an airport, a stop and a square.
\end{chapteropening}

\section[\texorpdfstring{\ru{прямо} (prýamo)}{прямо (prýamo)}]{\ru{прямо} (prýamo) — straight on}

\label{lesson:RU-C87-straight-on}



\begin{quote}
Before the new one: say the Russian for to the right, then the Russian for to the left.
\end{quote}

\subsection*{You'll want to know: \ru{прямо}}
\textbf{\ru{прямо}} (\emph{prýamo}) — \textquotedblleft{}straight on\textquotedblright{}.

\begin{cousinweb}
\textbf{\ru{прямо}} is \emph{\ru{прямой}}, \textquotedblleft{}straight\textquotedblright{}, made into an adverb. Vasmer calls the Proto-Slavic form disputed but thinks of a link with Greek \emph{prómos}, \textquotedblleft{}foremost\textquotedblright{}, and Old Norse \emph{framr}, \textquotedblleft{}forward\textquotedblright{}.
\end{cousinweb}

\begin{grammarlens}[title={What you've built}]
One more word for this chapter.
\end{grammarlens}

\subsection*{Your turn}
\begin{itemize}
\raggedright
  \item \emph{Say it:} \emph{prýamo}
  \item \emph{Say it:} \emph{prýamo}, once more
  \item \emph{Say it:} say \emph{prýamo}
  \item \emph{Recall:} say the Russian for to the right, then the Russian for to the left, then say \emph{prýamo} again
\end{itemize}

\begin{tcolorbox}[breakable,colback=teal!4,colframe=teal!35!black,title={Before you move on}]
What does \ru{прямо} mean? (\textquotedblleft{}Straight on\textquotedblright{}.) Say it once more.
\end{tcolorbox}
\section[\texorpdfstring{\ru{рядом} (ryádom)}{рядом (ryádom)}]{\ru{рядом} (ryádom) — nearby, next to}

\label{lesson:RU-C87-nearby}



\begin{quote}
Before the new one: say the Russian for to the left, then the Russian for straight on.
\end{quote}

\subsection*{You'll want to know: \ru{рядом}}
\textbf{\ru{рядом}} (\emph{ryádom}) — \textquotedblleft{}nearby, next to\textquotedblright{}.

\begin{cousinweb}
\textbf{\ru{рядом}} is \emph{\ru{ряд}}, \textquotedblleft{}a row\textquotedblright{}, in the form that says \textquotedblleft{}by, with\textquotedblright{}: \textquotedblleft{}in the same row\textquotedblright{}, side by side. Vasmer relates \emph{\ru{ряд}} to Lithuanian \emph{rinda}, \textquotedblleft{}a row, a line\textquotedblright{}.
\end{cousinweb}

\begin{grammarlens}[title={What you've built}]
One more word for this chapter.
\end{grammarlens}

\subsection*{Your turn}
\begin{itemize}
\raggedright
  \item \emph{Say it:} \emph{ryádom}
  \item \emph{Say it:} \emph{ryádom}, once more
  \item \emph{Say it:} say \emph{ryádom}
  \item \emph{Recall:} say the Russian for to the left, then the Russian for straight on, then say \emph{ryádom} again
\end{itemize}

\begin{tcolorbox}[breakable,colback=teal!4,colframe=teal!35!black,title={Before you move on}]
What does \ru{рядом} mean? (\textquotedblleft{}Nearby, next to\textquotedblright{}.) Say it once more.
\end{tcolorbox}
\section[\texorpdfstring{\ru{аэропорт} (aeropórt)}{аэропорт (aeropórt)}]{\ru{аэропорт} (aeropórt) — an airport}

\label{lesson:RU-C87-airport}



\begin{quote}
Before the new one: say the Russian for straight on, then the Russian for nearby.
\end{quote}

\subsection*{You'll want to know: \ru{аэропорт}}
\textbf{\ru{аэропорт}} (\emph{aeropórt}) — \textquotedblleft{}an airport\textquotedblright{}.

\begin{cousinweb}
\textbf{\ru{аэропорт}} comes from French \emph{aéroport}: \emph{aéro-}, from Greek \emph{a\'{\={e}}r}, \textquotedblleft{}air\textquotedblright{}, with \emph{port}, from Latin \emph{portus}, \textquotedblleft{}harbour\textquotedblright{}.
\end{cousinweb}

\begin{grammarlens}[title={What you've built}]
One more word for this chapter.
\end{grammarlens}

\subsection*{Your turn}
\begin{itemize}
\raggedright
  \item \emph{Say it:} \emph{aeropórt}
  \item \emph{Say it:} \emph{aeropórt}, once more
  \item \emph{Say it:} say \emph{aeropórt}
  \item \emph{Recall:} say the Russian for straight on, then the Russian for nearby, then say \emph{aeropórt} again
\end{itemize}

\begin{tcolorbox}[breakable,colback=teal!4,colframe=teal!35!black,title={Before you move on}]
What does \ru{аэропорт} mean? (\textquotedblleft{}An airport\textquotedblright{}.) Say it once more.
\end{tcolorbox}
\section[\texorpdfstring{\ru{остановка} (ostanóvka)}{остановка (ostanóvka)}]{\ru{остановка} (ostanóvka) — a stop (bus, tram)}

\label{lesson:RU-C87-stop}



\begin{quote}
Before the new one: say the Russian for nearby, then the Russian for an airport.
\end{quote}

\subsection*{You'll want to know: \ru{остановка}}
\textbf{\ru{остановка}} (\emph{ostanóvka}) — \textquotedblleft{}a stop (bus, tram)\textquotedblright{}.

\begin{cousinweb}
\textbf{\ru{остановка}} is \emph{\ru{остановить}}, \textquotedblleft{}to stop\textquotedblright{}, with the noun ending \emph{-\ru{ка}}. The verb is \emph{\ru{о}-} with \emph{-\ru{становить}}, from \emph{\ru{стать}}, \textquotedblleft{}to stand\textquotedblright{}, a cousin of Latin \emph{stāre}, \textquotedblleft{}to stand\textquotedblright{} (Vasmer).
\end{cousinweb}

\begin{grammarlens}[title={What you've built}]
One more word for this chapter.
\end{grammarlens}

\subsection*{Your turn}
\begin{itemize}
\raggedright
  \item \emph{Say it:} \emph{ostanóvka}
  \item \emph{Say it:} \emph{ostanóvka}, once more
  \item \emph{Say it:} say \emph{ostanóvka}
  \item \emph{Recall:} say the Russian for nearby, then the Russian for an airport, then say \emph{ostanóvka} again
\end{itemize}

\begin{tcolorbox}[breakable,colback=teal!4,colframe=teal!35!black,title={Before you move on}]
What does \ru{остановка} mean? (\textquotedblleft{}A stop (bus, tram)\textquotedblright{}.) Say it once more.
\end{tcolorbox}
\section[\texorpdfstring{\ru{площадь} (plóshchad)}{площадь (plóshchad)}]{\ru{площадь} (plóshchad) — a square}

\label{lesson:RU-C87-square}



\begin{quote}
Before the new one: say the Russian for an airport, then the Russian for a stop.
\end{quote}

\subsection*{You'll want to know: \ru{площадь}}
\textbf{\ru{площадь}} (\emph{plóshchad}) — \textquotedblleft{}a square\textquotedblright{}.

\begin{cousinweb}
\textbf{\ru{площадь}} — Old Russian and Church Slavonic \emph{\ru{площадь}}, from \emph{*ploskědĭ}, built on \emph{\ru{плоский}}, \textquotedblleft{}flat\textquotedblright{} (Vasmer): a square is a flat open place. The \emph{-\ru{щ}-} is the Church Slavonic sound.
\end{cousinweb}

\begin{grammarlens}[title={What you've built}]
One more word for this chapter.
\end{grammarlens}

\subsection*{Your turn}
\begin{itemize}
\raggedright
  \item \emph{Say it:} \emph{plóshchad}
  \item \emph{Say it:} \emph{plóshchad}, once more
  \item \emph{Say it:} say \emph{plóshchad}
  \item \emph{Recall:} say the Russian for an airport, then the Russian for a stop, then say \emph{plóshchad} again
\end{itemize}

\begin{tcolorbox}[breakable,colback=teal!4,colframe=teal!35!black,title={Before you move on}]
What does \ru{площадь} mean? (\textquotedblleft{}A square\textquotedblright{}.) Say it once more.
\end{tcolorbox}
